\section{Conclusion}

Surface-minimization models have proven remarkably successful at capturing local geometric motifs in a wide range of physical and biological networks. When formulated as smooth embedded manifolds minimizing an area functional subject to capacity or thickness constraints, these models explain branching angles, asymmetries, and stability regimes with a level of empirical accuracy that earlier wiring-length or one-dimensional optimization theories could not achieve.\cite{Meng2026SurfaceOptimization}

It is important to identify the limits of what surface minimization can decide in principle at the level of the classical variational formulation. When treated as a classical variational theory, surface minimization selects a single extremal configuration within a prescribed topology class and treats capacity constraints kinematically rather than resolving conflicts between competing continuation classes. As a result, questions of topology selection, permissibility of topology-changing transitions, and resolution of incompatible constraints lie outside the formal decision power of surface minimization alone, even under idealized conditions. This structural limitation is not unique to surface minimization, but is a general property of classical variational principles: optimization determines geometry within a prescribed configuration space, but does not determine the admissibility of the configuration space itself.

To resolve these indeterminacies, an additional decision layer is required. We have argued that this layer is best characterized as \emph{admissibility}: a principle governing which configurations and transitions are permitted under competing constraints. Admissibility is logically distinct from optimization and cannot be reduced to a refinement of the optimization functional. It operates categorically rather than incrementally, governs transitions rather than equilibria, and restricts the configuration space over which surface minimization is meaningfully applied. Introducing admissibility does not replace surface minimization; it completes it by specifying the domain within which optimization has determinate meaning.

This distinction leads to qualitative empirical consequences. In regimes where multiple continuation classes coexist, admissibility alters the space of expected behaviors, leading to testable signatures---including thresholded reorganization under continuous parameter variation, path dependence and hysteresis under reversal, and the persistence of locally non-optimal configurations due to forbidden transitions---that are not predictable from static surface minimization alone. Vascular anastomosis provides a canonical illustration of this contrast between geometric optimality and admissible reorganization, consistent with observations from recent three-dimensional imaging studies of vascular remodeling.\cite{Bentley2014} Experimentally, admissibility-governed behavior can be distinguished from pure optimization by protocols that slowly vary a control parameter (e.g., shear stress, growth demand, or pressure) while tracking topology in time: discontinuous or hysteretic topology changes at reproducible thresholds constitute signatures of admissibility, whereas purely optimization-driven geometry would vary continuously.

The admissibility concept articulated in this work is intentionally abstract, and this abstraction is a deliberate scope choice. Different physical systems may realize admissibility through distinct mechanisms---dynamical, developmental, energetic, stochastic, or evolutionary---and this paper does not attempt to unify or formalize those mechanisms. Its contribution is instead to clarify the common \emph{logical role} such mechanisms already play and to show that this role cannot be eliminated or absorbed into surface minimization without loss of explanatory power or predictive scope. In this sense, admissibility plays a role analogous to selection rules in other areas of physics: it determines which configurations or transitions are allowed to enter consideration, without prescribing their detailed dynamics. The present work concerns logical completeness rather than mechanism design; constructing system-specific admissibility criteria is deferred to future domain models. Accordingly, the novelty of the present work is not the observation that classical variational principles possess such structural limits---a point long recognized across physics---but the identification and formalization of this limit in the context of physical and biological network organization, where it has not previously been made explicit.

Thermal or stochastic fluctuations may, in practice, enable rare transitions between configuration classes. Such fluctuation-driven events do not eliminate the need for admissibility; rather, they provide one possible physical mechanism by which admissibility thresholds are crossed. Admissibility remains the logical criterion distinguishing permitted from forbidden transitions, while fluctuations determine how and how often those thresholds are reached.

We emphasize that this work does not propose admissibility as a new fundamental dynamical law, but rather identifies an unavoidable logical role that existing physical, developmental, or material mechanisms already implement in practice.

Although the discussion has focused primarily on biological networks, the admissibility distinction applies equally to other physical systems modeled by constrained minimal surfaces, including material fracture patterns, fluidic transport networks, and engineered microfluidic architectures, wherever topology-changing transitions arise under competing constraints.

Surface minimization and admissibility therefore occupy complementary positions in the theoretical hierarchy. Surface minimization determines geometry within admissible classes, while admissibility determines which classes and transitions are physically available. Both principles are required for a complete theory of physical network organization, and neither subsumes the other. Future work may operationalize admissibility in specific systems, bridging it with dynamical or stochastic models to predict reorganization thresholds and transition boundaries. Formally, surface minimization determines geometry conditional on admissibility, whereas admissibility itself is not derivable from the variational principle and must be specified independently.
